% SPDX-License-Identifier: AGPL-3.0-or-later
% Commercial license available
% © Concepts 1996–2026 Miroslav Šotek. All rights reserved.
% © Code 2020–2026 Miroslav Šotek. All rights reserved.
% ORCID: 0009-0009-3560-0851
% Contact: www.anulum.li | protoscience@anulum.li
% SCPN-FUSION-CORE — regime-conditional conformal coverage of a transport surrogate.
\documentclass[11pt]{article}

\usepackage[utf8]{inputenc}
\usepackage[T1]{fontenc}
\usepackage{amsmath}
\usepackage{booktabs}
\usepackage{microtype}
\usepackage{natbib}
\usepackage[colorlinks=true,linkcolor=blue,citecolor=blue,urlcolor=blue]{hyperref}
\usepackage[margin=1in]{geometry}
\bibliographystyle{unsrtnat}

\title{A Marginal Conformal Certificate for a Neural Transport Surrogate\\
Hides a False-Stable Bias:\\
Regime-Conditional Coverage and Class-Conditional Stability Sets}
\author{Miroslav \v{S}otek\\
\small ANULUM CH\&LI, Marbach SG, Switzerland, \url{https://www.anulum.li}\\
\small ORCID: \href{https://orcid.org/0009-0009-3560-0851}{0009-0009-3560-0851}\\
\small \href{mailto:protoscience@anulum.li}{protoscience@anulum.li}}
\date{Preprint --- 16 November 2026}

\begin{document}
\maketitle

\begin{abstract}
A split-conformal certificate gives a neural surrogate of turbulent transport
fluxes, trained on the QLKNN-10D database, 90\% marginal coverage per channel:
0.901, 0.900 and 0.903 on 47,828 held-out rows for $\chi_e$, $\chi_i$ and
$D_e$.  We ask whether that coverage holds across transport regimes.  It does
not.  Rows with zero flux are covered at 0.995--0.997, while rows with nonzero
flux are covered at only 0.697--0.738, and those above a near-threshold band at
0.665--0.711.  Three alternatives to the certificate, declared in advance, do not repair
this: Mondrian groups by an input-gradient label, and scores normalised by the
gyro-Bohm scale or by the prediction.  Two further variants, chosen after
seeing those results and reported as such, do not repair it either.  The
cause is the point model: it predicts a flux no larger than the certificate
radius on 64.5--73.8\% of the rows above the band, a false-stable bias that no
score built from the prediction can separate from true stability.
Class-conditional conformal stability sets, which use the prediction only to
rank rows, cover each true class at 0.891--0.903.  Flux intervals conditioned
on those sets cover the nonzero-flux rows at 0.792--0.797, against
0.697--0.738 for the certificate on the same rows.  The cost is width: the
median interval on those rows is 12.2, 14.4 and 2.36\,m$^2$\,s$^{-1}$.  Over 20
random re-splits the per-class coverage averages 0.900--0.901.  Propagated
with frozen coefficients through one L-mode transport case, the certificate
interval spans a core temperature of 0.262--7.224\,keV around 1.702\,keV, so
it carries little information as a profile prediction.  All guarantees are
marginal within each class and assume exchangeability.
\end{abstract}

\noindent\textbf{Keywords:} conformal prediction; uncertainty quantification;
neural surrogate; turbulent transport; QLKNN; label-conditional validity

\section{Introduction}

Neural surrogates of quasilinear transport models replace the transport code
inside integrated modelling at a fraction of the cost.  QLKNN fits the
QuaLiKiz model with networks whose training is shaped by the physics of the
critical gradient \citep{vandePlassche2020}.  Its cost function fits the fluxes
only on unstable points and adds an explicit penalty only for positive
predictions where QuaLiKiz is stable (their Appendix~E).  False-stable
predictions are penalised only through the squared error on unstable points,
and no term controls their rate.  Recent work estimates the uncertainty of
such surrogates with model-based methods, for example a classifier for mode
presence and a Bayesian regressor for each flux \citep{Ho2025}.

Split conformal prediction gives distribution-free intervals around any
trained model.  The guarantee is marginal, not conditional on the input, and
it assumes that calibration and test rows are exchangeable
\citep{Gopakumar2026}.  For a transport surrogate the regimes matter: in the
database used here most points have zero flux, and an interval that is right
on average can be wrong where the flux is not zero.  Precise conditional
validity on the input is not achievable in a useful way, but validity
conditional on the label is: a label-conditional inductive conformal predictor
controls the set-valued analogues of the false-positive and false-negative
rates \citep{Vovk2012}.

We measure the coverage of an existing 90\% certificate per regime, try the
standard remedies, identify why they fail, and give a label-conditional
remedy with its cost.  The surrogate is not QLKNN itself, and nothing here is
attributed to QLKNN.

\section{Surrogate, data and certificate}

The surrogate is a fully connected network with hidden layers of 1024, 512 and
256 units and three outputs, $\chi_e$, $\chi_i$ and $D_e$, trained on the
QLKNN-10D database \citep{QLKNN10D2019}.  It has no gate and no log transform,
and it converts gyro-Bohm-normalised outputs to physical units.  Its metrics
record gives relative $L_2$ errors of 0.0917, 0.0954 and 0.0943 on the
training, validation and test splits.  The cost function that produced these
weights is not recorded with them.  The certificate uses the processed validation and
test splits, of 47,826 and 47,828 rows; it records each file's SHA-256 digest.  Negative fluxes were clipped to zero in preprocessing, so the
zero-flux class includes clipped rows.

The certificate uses the absolute residual $|y - \hat{y}|$ as the score on the
47,826 validation rows.  The radius $r$ of each channel is the calibration
score of rank $\lceil (n+1)(1-\alpha) \rceil = 43{,}045$ for $\alpha = 0.1$, and
the interval is $[\max(0, \hat{y} - r), \hat{y} + r]$.  The radii are 2.92,
3.30 and 0.394\,m$^2$\,s$^{-1}$.  The 47,828 test rows serve only to check
coverage.

We use three target-defined regimes per channel: zero flux; a near-threshold
band $(0, q]$, with $q$ the 10th percentile of the positive calibration
targets; and the rows above the band.  These regimes are not known at
prediction time.  They serve for evaluation only, never for calibration.  The
database also carries an input-gradient label (ITG for $R/L_{T_i} > 4$, TEM for
$R/L_{T_e} > 5$ with $R/L_{T_i} \le 4$, neither otherwise).  That label is
known at prediction time, but it is not a stability label.

\section{Coverage per regime}

Table~\ref{tab:regimes} lists the coverage on the test rows.  The certificate
reproduces its marginal coverage exactly.  The constant radius is set mainly by
the zero-flux rows, 63--68\% of the data, and it over-covers them.  It
falls short of 0.90 on the rows with nonzero flux by 16.2 to 20.3 points, and
on the rows above the band by 18.9 to 23.5 points.  In the near-threshold band the intervals are so
wide relative to the target that they carry no information: their median
relative width $2r/y$ is 216, 12,214 and 296.  Under the input labels the
coverage ranges from 0.862 to 0.981.

\begin{table}[ht]
\centering
\caption{Coverage of the 90\% certificate on the 47,828 test rows, with the
row count in parentheses.  Nonzero: the near-threshold band and the rows
above it together.  Labels: the range over the three input-gradient labels.}
\label{tab:regimes}
\small
\begin{tabular}{lrrrrrr}
\toprule
Channel & Marginal & Zero flux & Near threshold & Above band & Nonzero & Labels \\
\midrule
$\chi_e$ & 0.901 & 0.997 (30,105) & 0.999 (1,748) & 0.709 (15,975) & 0.737 (17,723) & 0.862--0.969 \\
$\chi_i$ & 0.900 & 0.995 (30,105) & 0.991 (1,763) & 0.711 (15,960) & 0.738 (17,723) & 0.862--0.981 \\
$D_e$ & 0.903 & 0.997 (32,758) & 0.993 (1,452) & 0.665 (13,618) & 0.697 (15,070) & 0.882--0.956 \\
\bottomrule
\end{tabular}
\end{table}

\section{Standard conformal variants}

Four variants were declared in the code before any result was seen: the
marginal absolute score (the certificate); Mondrian groups by the input label;
the absolute score divided by the gyro-Bohm scale; and the absolute score
divided by $|\hat{y}| + r_0$, with $r_0$ the certificate radius.  After seeing
their results we added two Mondrian variants grouped by the predicted flux
($\hat{y} \le r_0$ or $\hat{y} > r_0$), one with each score.  They are marked
post hoc in the code, in the report and here.

\begin{table}[ht]
\centering
\caption{Conformal variants at the 90\% level.  Coverage of the rows above the
near-threshold band and of all nonzero-flux rows, and the median half-width in
m$^2$\,s$^{-1}$, for $\chi_e$ / $\chi_i$ / $D_e$.  The last two rows are post
hoc.}
\label{tab:variants}
\small
\begin{tabular}{lccc}
\toprule
Variant & Above band & Nonzero & Median half-width \\
\midrule
Marginal absolute (certificate) & 0.709 / 0.711 / 0.665 & 0.737 / 0.738 / 0.697 & 2.92 / 3.30 / 0.394 \\
Mondrian, input label & 0.713 / 0.716 / 0.671 & 0.741 / 0.743 / 0.702 & 3.28 / 5.73 / 0.548 \\
Gyro-Bohm normalised & 0.701 / 0.709 / 0.659 & 0.730 / 0.737 / 0.691 & 2.46 / 2.68 / 0.326 \\
$|\hat{y}| + r_0$ normalised & 0.706 / 0.723 / 0.664 & 0.735 / 0.749 / 0.696 & 1.24 / 1.14 / 0.196 \\
Mondrian, predicted flux & 0.724 / 0.733 / 0.669 & 0.751 / 0.758 / 0.700 & 0.893 / 0.658 / 0.131 \\
\quad with $|\hat{y}| + r_0$ score & 0.719 / 0.735 / 0.666 & 0.747 / 0.760 / 0.698 & 0.716 / 0.541 / 0.110 \\
\bottomrule
\end{tabular}
\end{table}

No variant lifts the coverage above the band beyond 0.735
(Table~\ref{tab:variants}).  The $|\hat{y}| + r_0$ score is the efficient
choice: it cuts the median width by half or more at about the same coverage.  The reason none of
them helps lies in the point model.  On the rows above the band the surrogate
predicts at most $r_0$ for 68.3\%, 64.5\% and 73.8\% of them, with a median
ratio of prediction to target of 0.493, 0.602 and 0.305.  Such rows look
stable to any score or group built from the prediction, so no such score or
group can separate them from the truly stable rows.

\section{Class-conditional stability sets}

The prediction still ranks the rows even where its magnitude is wrong.  We
therefore build, per channel, a label-conditional conformal predictor
\citep{Vovk2012} for the stability class.  The score is $\hat{y}$ for the label
``stable'' and $-\hat{y}$ for ``unstable''.  For each class $c$ the threshold
$q_c$ is the calibration score of rank $\lceil (n_c+1)(1-\alpha) \rceil$ among
the calibration rows of that class.  A label enters a test row's set when the
row's score does not exceed the threshold of that label.  Under exchangeability
the set contains the true label with probability at least $1-\alpha$ in each
class.  The flux interval then follows the set:
\begin{itemize}
\item $[0, 0]$ when the set is \{stable\};
\item $[\max(0, \hat{y} - q_f), \hat{y} + q_f]$ when it is \{unstable\};
\item $[0, \hat{y} + q_f]$ when it holds both labels;
\item an abstention when it is empty.
\end{itemize}
Here $q_f$ is the absolute-residual quantile of the unstable calibration rows.
A union bound then covers an unstable row's flux with probability at least
$1-2\alpha = 0.80$.  This experiment was designed after the variant comparison
had shown the bias; its settings were fixed before it was run.

\begin{table}[ht]
\centering
\caption{Stability sets on the test rows (original split).  Set coverage and
the share of rows whose set holds both labels, per true class.  Flux coverage
of the nonzero-flux rows, with the certificate on the same rows for
comparison, and the median full width of their intervals in
m$^2$\,s$^{-1}$.  No set was empty.}
\label{tab:sets}
\small
\begin{tabular}{lcccccc}
\toprule
 & \multicolumn{2}{c}{Set coverage} & \multicolumn{2}{c}{Both labels} & \multicolumn{2}{c}{Nonzero-flux rows} \\
\cmidrule(lr){2-3}\cmidrule(lr){4-5}\cmidrule(lr){6-7}
Channel & stable & unstable & stable & unstable & flux coverage (cert.) & median width \\
\midrule
$\chi_e$ & 0.901 & 0.900 & 0.113 & 0.101 & 0.797 (0.737) & 12.2 \\
$\chi_i$ & 0.902 & 0.894 & 0.205 & 0.156 & 0.792 (0.738) & 14.4 \\
$D_e$ & 0.903 & 0.891 & 0.081 & 0.071 & 0.793 (0.697) & 2.36 \\
\bottomrule
\end{tabular}
\end{table}

Table~\ref{tab:sets} shows the result.  The flux coverage of the nonzero-flux
rows rises to 0.792--0.797, close to the 0.80 bound, from 0.737, 0.738 and
0.697: by 5 to 10 points.
The price is width: the unstable residual quantile $q_f$ is 11.6, 13.8 and
2.32\,m$^2$\,s$^{-1}$, against certificate radii of 2.92, 3.30 and 0.394.  Two
unstable set coverages of the original split lie below 0.90 (0.894 and 0.891).
Over 20 random halvings of the pooled calibration and test rows, the
per-class set coverage has a mean of 0.900--0.901 and a range of 0.892--0.909,
and the nonzero-flux coverage has a mean of 0.800--0.803 and a range of
0.791--0.810.  The guarantee therefore holds on average; the original split's
$D_e$ value of 0.891 lies just below the re-split minimum of 0.894.

\section{Propagation to a transport profile}

To see what the intervals mean for a profile, we ran one L-mode case of the
integrated transport solver: the ITER-like test configuration at 20\,MW of
auxiliary power, with neoclassical diffusivity and 1200 steps of 10\,ms.  The
power lies below the threshold where the solver's pedestal modifier acts.  The
neural closure reaches a steady state with a core temperature of 1.702\,keV, a
mean temperature of 0.957\,keV and an energy confinement time of 1.242\,s.
We then froze the turbulent diffusivity at its lower and upper interval bounds
and ran the solver to steady state with those fixed coefficients.  The
neoclassical part stays unchanged.  The lower bound is clipped at zero.

At the certificate's bounds the core temperature spans 0.262--7.224\,keV and
the confinement time 0.246--4.533\,s.  The certificate radii exceed the median
turbulent diffusivities of the case (0.181, 0.368 and 0.0152\,m$^2$\,s$^{-1}$)
many times over.  The lower ends of both intervals give the same steady state,
consistent with the turbulent part being clipped to zero.  The narrower
$|\hat{y}| + r_0$ interval moves only the upper end, to 0.455\,keV.  As a
profile prediction, a 90\% flux interval of this kind spans a factor of 4.2
upward and 6.5 downward in core temperature.  This is a frozen-coefficient sensitivity
for one case, not a self-consistent propagation: the diffusivity depends on the
evolving gradients.

\section{Limitations}

\begin{itemize}
\item All guarantees are marginal within a class and assume that calibration
and test rows are exchangeable.  Nothing is claimed for data outside the
database's distribution.
\item The targets are QuaLiKiz-derived database values, not gyrokinetic or
experimental fluxes.
\item The surrogate is one network.  Its training cost is not recorded with
its weights, so the bias is not attributed to any training choice.
\item The two post hoc variants and the stability sets were designed after
earlier results.  Their settings were fixed before their own runs.
\item The profile propagation covers one case, with frozen coefficients.
\end{itemize}

\section*{Software version and date}

This preprint describes \texttt{scpn-fusion} at commit \texttt{8191a1d31e39}
of 16 November 2026, on the main line after release 4.1.0.  The analysis
modules are not part of that release.  Each report records the commit that
produced it, and the evidence manifest lists those commits.  Since then the
analysis code changed only by making private helpers public.  The variant
report, produced after the first such change, reproduces the regime report's
certificate coverage exactly.

\section*{Data and code availability}

Code is available at \url{https://github.com/anulum/scpn-fusion-core}.  The
certificate and the four reports cited here, with their source revisions, Git
blob identities and SHA-256 digests, are in \texttt{evidence/} of this
submission directory.  The generator in that directory checks every quoted
value against them, including the derived nonzero-flux coverages.  This is a
venue-neutral preprint and has not been submitted.  No DOI has been reserved
and the manuscript-content licence has not yet been selected by the author.

\section*{Use of AI tools}

Large language model coding agents were used substantively for generating and
debugging the analysis code and for drafting and revising the manuscript text.
The agents were commercially available models, which are not named because no
vendor is involved in or sponsors this work.  They were coordinated through
SYNAPSE CHANNEL (version 0.99.52 at the time of writing;
\url{https://github.com/anulum/synapse-channel}), the open-source
agent-coordination tool of the ANULUM Agentic Coordination, Assurance and
Continuity Systems group.  The author specified the analyses, verified the
code through its tests and every reported number against the committed
reports, and takes full responsibility for the content.

\bibliography{references}
\end{document}
